\documentclass[12pt]{article}

% Layout.
\usepackage[top=1.2in, bottom=0.75in, left=1in, right=1in, headheight=1.0in, headsep=0pt]{geometry}

% Fonts.
\usepackage{mathptmx}
\usepackage[scaled=0.86]{helvet}
\renewcommand{\emph}[1]{\textsf{\textbf{#1}}}

% TiKZ.
\usepackage{tikz, pgfplots}
\usetikzlibrary{calc}
\pgfplotsset{my style/.append style={axis x line=middle, axis y line=middle, xlabel={$x$}, ylabel={$y$}}}
\pgfplotsset{compat=1.16}

% Misc packages.
\usepackage{amsmath,amssymb,latexsym}
\usepackage{graphicx}
\usepackage{array}
\usepackage{xcolor}
\usepackage{multicol}

% Commands to set various header/footer components.
\makeatletter
\def\doctitle#1{\gdef\@doctitle{#1}}
\doctitle{Use {\tt\textbackslash doctitle\{MY LABEL\}}.}
\def\docdate#1{\gdef\@docdate{#1}}
\docdate{Use {\tt\textbackslash docdate\{MY DATE\}}.}
\def\doccourse#1{\gdef\@doccourse{#1}}
\let\@doccourse\@empty
\def\docscoring#1{\gdef\@docscoring{#1}}
\let\@docscoring\@empty
\def\docversion#1{\gdef\@docversion{#1}}
\let\@docversion\@empty
\makeatother

% Headers and footers layout.
\makeatletter
\usepackage{fancyhdr}
\pagestyle{fancy}
\fancyhf{} % Clears all headers/footers.
\lhead{\emph{\@doctitle\hfill\@docdate} \medskip
\ifnum \value{page} > 1\relax\else\\
\emph{Name: \rule{3.5in}{0.8pt}\hfill \@docscoring}
\fi}

\rfoot{\emph{\@docversion}}
\lfoot{\emph{\@doccourse}}
\cfoot{\emph{\thepage}}
\renewcommand{\headrulewidth}{0pt}%
\makeatother

% Paragraph spacing
\parindent 0pt
\parskip 6pt plus 1pt

% A problem is a section-like command. Use \problem{5} for a problem worth 5 points.
\newcounter{probcount}
\newcounter{subprobcount}
\setcounter{probcount}{0}

\newcommand{\problem}[1]{%
\par
\addvspace{4pt}%
\setcounter{subprobcount}{0}%
\stepcounter{probcount}%
\makebox[0pt][r]{\emph{\arabic{probcount}.}\hskip1ex}\emph{[#1 points]}\hskip1ex}

\newcommand{\thesubproblem}{\emph{\alph{subprobcount}.}}

% like \problem but with name
\newcommand{\nameprob}[2]{%
\par
\addvspace{4pt}%
\setcounter{subprobcount}{0}%
\stepcounter{probcount}%
\makebox[0pt][r]{\emph{#1.}\hskip1ex}\emph{[#2 points]}\hskip1ex}

% Subproblems are an enumerate-like environment with a consistent
% numbering scheme.  Use \begin{subproblems}\item...\item...\end{subproblems}
\newenvironment{subproblems}{%
\begin{enumerate}%
\setcounter{enumi}{\value{subprobcount}}%
\renewcommand{\theenumi}{\emph{\alph{enumi}}}}%
{\setcounter{subprobcount}{\value{enumi}}\end{enumerate}}

% Blanks for answers in normal and math mode.
\newcommand{\blank}[1]{\rule{#1}{0.75pt}}
\newcommand{\mblank}[1]{\underline{\hspace{#1}}}
\def\emptybox(#1,#2){\framebox{\parbox[c][#2]{#1}{\rule{0pt}{0pt}}}}

% Misc.
\renewcommand{\d}{\displaystyle}
\newcommand{\ds}{\displaystyle}


\doctitle{Math 252: Quiz 5}
\docdate{12 February 2026}
\doccourse{}
\docversion{}
\docscoring{\fbox{{\LARGE \strut}\blank{0.8in} / 25}}

\begin{document}
\phantom{foo} \vspace{-5mm}

30 minutes maximum.  No aids (internet, book, etc.) are permitted, but \textbf{please see the trigonometric identities on the final page!}  For full credit please show all work, use proper notation, and put answers in reasonably-simplified form.  25 points possible.

\problem{15}  {\large Evaluate the following integrals.}

\begin{subproblems}
\item  {\large $\ds \int x\, e^{2x}\,dx = $}
\vfill

\item  {\large $\ds \int \cos^3 w\, \sin^4 w\,dw = $}
\vfill

\item  {\large $\ds \int_0^1 \arctan x\,dx = $}
\vfill

\clearpage\newpage
\item  {\large $\ds \int e^x\, \sin x\,dx = $}
\vspace{65mm}

\item  {\large $\ds \int \tan^2 x\,dx = $}
\vfill
\end{subproblems}

\problem{4}  {\large We say two functions $f(x),g(x)$ are \textsl{orthogonal} on the interval $[-\pi,\pi]$ if the integral of their product is zero: $\int_{-\pi}^\pi f(x)\,g(x)\,dx = 0$.  Show that the functions $\sin(2x)$ and $\cos(3x)$ are orthogonal on the interval $[-\pi,\pi]$.}
\vspace{90mm}

\newpage
\problem{6}  {\large \strut Sketch the region between \,$y=\cos x$\, and the $x$-axis on the interval $0\le x \le \pi/2$.  Find the volume of the solid which results by rotating the region around the $x$-axis.  (\textsl{Hint.  Use disks.})}
\vfill

\newpage
\nameprob{EC}{1} (\emph{Extra Credit})  Assume $n$ is a large and positive integer.  One of these indefinite integrals is much easier than the other.  \emph{Circle the \underline{easier} one, and then evaluate it.}
    $$\int \tan^n x\, \sec x\,dx \hspace{1.5in} \int \sec^n x\, \tan x\,dx$$
\vspace{2.5in}

\noindent \hrule

\medskip
\noindent {\small You may find the following \textbf{trigonometric formulas} useful.  Other formulas, not listed here, should be in your memory, or you can derive them from the ones here.
\begin{align*}
\sin(\alpha \pm \beta) &= \sin \alpha \cos \beta \pm \cos \alpha \sin \beta \qquad &
    \sin(ax) \sin(bx) &= \frac{1}{2} \cos((a-b)x) - \frac{1}{2} \cos((a+b)x) \\
\cos(\alpha \pm \beta) &= \cos \alpha \cos \beta \mp \sin \alpha \sin \beta \qquad &
    \sin(ax) \cos(bx) &= \frac{1}{2} \sin((a-b)x) + \frac{1}{2} \sin((a+b)x) \\
 && \cos(ax) \cos(bx) &= \frac{1}{2} \cos((a-b)x) + \frac{1}{2} \cos((a+b)x)
\end{align*}
}

\noindent \hrule

\bigskip
\centerline{\footnotesize \textsc{blank space}}
\vfill
\end{document}